% ---------------------------------------------------------------------
\section{重到轻形状因子与无粲 $B$ 衰变的 QCD 因子化}\label{sec:heavy-qcdf}

%%INSERT fig_qcdf.tex

两个末态介子都很轻（$\bar B\to\pi\pi$、$\pi K$ 等）时，旁观反夸克进入高能轻介子 $M_1$ 须交换虚度 $\sim m_b\Lqcd$ 的胶子，出现硬共线标度与同阶的硬旁观者散射（hard spectator scattering，\cref{fig:heavy-qcdf}）。

\begin{proposition}[重到轻形状因子：软重叠与硬旁观者 \tagA]\label{prop:heavy-btopi-ff}
设 $\bar B\to\pi$ 跃迁中 $\pi$ 在 $B$ 静止系中的能量 $E=\frac{m_B^2-q^2}{2m_B}\sim m_b/2\gg\Lqcd$（大反冲）。则在 $\Lqcd/m_b$ 的领头幂次、对 $\as$ 的所有阶，每个形状因子 $F_i$（$i=+,0,T$）可写成
\begin{equation}\label{eq:heavy-btopi-ff}
F_i^{B\to\pi}(E)=C_i(E,\mu)\,\xi_\pi(E,\mu)+\int_0^\infty\frac{\dd\omega}{\omega}\int_0^1\dd u\;T_i(E,\omega,u,\mu)\,\phi_B^+(\omega,\mu)\,\phi_\pi(u,\mu),
\end{equation}
其中 $\xi_\pi$ 是与 $i$ 无关的普适“软重叠”（soft overlap）形状因子，在 SCET 中定义为包含硬共线、共线与软效应的矩阵元；$C_i$ 是硬系数（含与 $i$ 有关的运动学因子）；$T_i=\order{\as}$ 是硬系数与硬共线喷注函数的卷积；$\phi_B^+$ 是 $B$ 介子的领头光锥分布振幅（$\omega$ 为旁观者光锥动量，$\int_0^\infty\frac{\dd\omega}{\omega}\phi_B^+=\lambda_B^{-1}$\cite{Grozin:1996pq}，演化见\cite{Lange:2003ff}），$\phi_\pi$ 见\cref{def:excl-da}。两项都 $\sim(\Lqcd/m_b)^{3/2}$。
非微扰地位：\cref{asm:intro-bridge}（MA1、MA2）。
% errtype: E-S
\uses{def:excl-da, thm:scet-decoupling, asm:intro-bridge}
\end{proposition}

\begin{proofsketch}
(1) 两步匹配 QCD $\to$ SCET$_{\rm I}$ $\to$ SCET$_{\rm II}$；在 SCET$_{\rm I}$ 中 $B$ 的软组分相对于硬共线模式是超软的，经 BPS 场重定义退耦（\cref{thm:scet-decoupling}）。(2) 只含一个共线夸克场 $\bar\xi W$ 的重到轻流，因 $\sla vh_v=h_v$、$\sla n_A\xi=0$ 的约束，对各形状因子只差运动学因子，给出 $\xi_\pi$；多含一个硬共线胶子场的流在积出硬共线模式后化为喷注函数与 $\phi_B^+$、$\phi_\pi$（\cref{def:excl-da}）的卷积。(3) 前一类矩阵元若再写成卷积则在端点发散；Beneke 与 Feldmann 把这些效应放进 $\xi_\pi$ 的定义，得到对 $\as$ 所有阶成立的\cref{eq:heavy-btopi-ff}\cite{Beneke:2003pa}；BPS 用另一种标度分离得到无端点奇异性的硬核\cite{Bauer:2002aj}（E-S 型陈述；另见\cite{Hill:2002vw,Beneke:2005gs}）。
\end{proofsketch}

\paragraph{地位与缺口}
全阶论证（[A]），前提是\cref{rem:scet-assumptions}的 SCET 前提（模式完备性，E-S 型陈述）。被因子化的只是“软重叠 + 硬旁观者”的分解：$\xi_\pi$ 不能再因子化为分布振幅的卷积（端点发散；同类现象见\cref{prop:heavy-qcdf-endpoint}与\cref{prop:excl-endpoint}，[P]）；对 $B_c\to\eta_c$，软重叠的全阶双对数修正已可在微扰的非相对论框架中分析\cite{Bell:2024bxg}。升级途径：对 $B\to\pi$ 证明模式完备性。

\paragraph{记号与幂次计数。}
以 $b\to s$ 为例，$\mathcal H_{\rm eff}=\frac{G_F}{\sqrt2}\sum_{p}\lambda_p(C_1Q_1^p+C_2Q_2^p+\cdots)$，$Q_1^p=(\bar pb)_{V-A}(\bar sp)_{V-A}$、$Q_2^p=(\bar p_ib_j)_{V-A}(\bar s_jp_i)_{V-A}$，NDR 方案 $\mu=m_b$ 处 $C_1=1.081$、$C_2=-0.190$\cite[表~1]{BenekeEtAl2001}；$\Phi_M(x)=6x\bar x[1+\sum_n\alpha_n^MC_n^{(3/2)}(2x-1)]$，扭度 3 振幅 $\Phi_p=1$，$r_\chi^M=2\mu_M/m_b$，$\int_0^1\frac{\dd\xi}{\xi}\Phi_B=\frac{m_B}{\lambda_B}$；$M_1$ 接收旁观者，$M_2$ 为发射的介子。由 $F^{B\to\pi}(0)\sim(\Lqcd/m_b)^{3/2}$、$f_B\sim\Lqcd^{3/2}m_b^{-1/2}$、$f_\pi\sim\lambda_B\sim\Lqcd$，形状因子项与硬旁观者项同为 $\Lqcd^{5/2}m_b^{1/2}$。旁观者（动量 $\xi p_B$，$\xi\sim\Lqcd/m_b$）以分数 $\bar y=1-y$ 进入 $M_1$ 时，胶子动量 $k=\xi p_B-\bar yp_1$ 满足
\begin{equation}\label{eq:heavy-hc-gluon}
k^2=\xi^2m_B^2-\bar y\,\xi\,m_B^2\sim-m_b\Lqcd\qquad(\bar y=\order{1}),
\end{equation}
是硬共线的，BBNS 因此在 $\mu_h=\sqrt{\Lambda_h\mu}$（$\Lambda_h=0.5$\,GeV）处取其 $\as$\cite{BenekeEtAl2001}。

\begin{proposition}[无粲两体衰变的领头幂次 QCD 因子化 \tagP]\label{prop:heavy-qcdf-lp}
设 $M_1$、$M_2$ 为轻介子。在有限阶的意义上，BBNS 因子化公式在 $\Lqcd/m_b$ 的领头幂次成立：
\begin{equation}\label{eq:heavy-qcdf}
\begin{split}
\langle M_1M_2|Q_i|\bar B\rangle={}&F^{B\to M_1}(0)\,f_{M_2}\!\int_0^1\!\dd x\,T^{\rm I}_i(x)\Phi_{M_2}(x)\\
&+f_Bf_{M_1}f_{M_2}\!\int_0^1\!\dd\xi\,\dd x\,\dd y\;T^{\rm II}_i(\xi,x,y)\,\Phi_B(\xi)\Phi_{M_1}(y)\Phi_{M_2}(x)
\end{split}
\end{equation}
（旁观者可进入任一介子时再加 $M_1\leftrightarrow M_2$ 的形状因子项），$T^{\rm I}=1+\order{\as}$，$T^{\rm II}=\order{\as}$。具体地：(a) 在 $\order{\as}$，“非因子化”图的软与共线发散在领头幂次相消，顶点修正给出
\begin{gather}
a_{1,\rm I}=C_1+\frac{C_2}{N_c}\Big[1+\frac{C_F\as}{4\pi}V_{M_2}\Big],\qquad a_{2,\rm I}=C_2+\frac{C_1}{N_c}\Big[1+\frac{C_F\as}{4\pi}V_{M_2}\Big],\label{eq:heavy-ai}\\
V_M=12\ln\frac{m_b}{\mu}-18+\int_0^1\dd x\,g(x)\,\Phi_M(x)\quad(\text{NDR}),\label{eq:heavy-VM}\\
\int_0^1\dd x\,g(x)\Phi_M(x)=-\frac12-3\ii\pi+\Big(\frac{11}{2}-3\ii\pi\Big)\alpha_1^M-\frac{21}{20}\alpha_2^M+\cdots,\label{eq:heavy-gegen}
\end{gather}
% VF-BQ-1, VF-BQ-2
$g(x)=3\big(\frac{1-2x}{1-x}\ln x-\ii\pi\big)+\big[2\,\mathrm{Li}_2(x)-\ln^2x+\frac{2\ln x}{1-x}-(3+2\ii\pi)\ln x-(x\leftrightarrow1-x)\big]$（约定同\cite{BenekeEtAl2001}），系数 $12$ 由 $a_{1,2}$ 与 $\mu$ 无关唯一确定（对任意 $N_c$）；(b) $T^{\rm I}$ 的红外有限性已显式检验到 NNLO（两圈），$T^{\rm II}$ 的到单圈（$\order{\as^2}$）。
非微扰地位：\cref{asm:intro-bridge}（MA1、MA2）。
% errtype: E-S
\uses{prop:heavy-btopi-ff, def:excl-da, asm:intro-bridge}
\end{proposition}

\begin{proofsketch}
(a) 是 BBNS 的单圈计算\cite[式~(46)--(48)]{BenekeEtAl2001}，$F^{B\to M_1}$ 按\cref{prop:heavy-btopi-ff}理解，$\Phi_M$ 见\cref{def:excl-da}。系数 $12$ 可独立导出：由 $\frac{\dd\bm C}{\dd\ln\mu}=\frac{\as}{4\pi}\gamma^{\tp}\bm C$，$\gamma=\begin{psmallmatrix}-6/N_c&6\\6&-6/N_c\end{psmallmatrix}$（$C_1\pm C_2$ 的本征值 $\pm6(N_c\mp1)/N_c$），
\begin{equation}\label{eq:heavy-a1-rg}
\frac{\dd}{\dd\ln\mu}\Big(C_1+\frac{C_2}{N_c}\Big)=\frac{\as}{4\pi}\Big[\Big(-\frac{6}{N_c}C_1+6C_2\Big)+\frac1{N_c}\Big(6C_1-\frac6{N_c}C_2\Big)\Big]=\frac{\as}{4\pi}\,\frac{12C_F}{N_c}\,C_2 ,
\end{equation}
恰被 $\frac{C_2}{N_c}\frac{C_F\as}{4\pi}12\ln\frac{m_b}{\mu}$ 的导数抵消；$a_{2,\rm I}$ 同理，常数 $-18$ 依赖方案；Gegenbauer 项以 $10^{-20}$ 精度重现（脚本 \texttt{s2\_\allowbreak qcdf.py}）。(b) 见\cite{Bell:2007tv,Bell:2009nk,Beneke:2009ek,Bell:2015koa,Bell:2020qus}（两圈顶点与企鹅）、\cite{Beneke:2005vv,Beneke:2006mk,Kivel:2006xc,Pilipp:2007mg,Beneke:2005gs}（旁观者散射）。
\end{proofsketch}

在 $\mu=m_b$ 取渐近分布振幅，$V_\pi=-18.5-3\ii\pi$：领头幂次的强相位是 $\order{\as}$ 的微扰量；以 $\as(4.2\,{\rm GeV})=0.223$ 得 $a_{2,\rm I}\approx0.012-0.080\ii$，色压低的 $a_2$ 中修正与领头项同量级。

\paragraph{地位与缺口}
部分（[P]）：在 $\order{\as}$ 阶证明了 BBNS 公式，$T^{\rm I}$ 到 NNLO、$T^{\rm II}$ 到 $\order{\as^2}$ 显式红外有限，但对两个轻介子没有公认的全阶证明（BBNS 本人指出须把 $B\to D\pi$ 的两圈证明推广到这一情形\cite[第~8~节]{BenekeEtAl2000}）。SCET 中有领头幂次的全阶推导：Chay 与 Kim 论证振幅对 $\as$ 的所有阶因子化\cite{Chay:2003ju}；BPRS 认为其定理在 $\as(E_M\Lqcd)$ 不可微扰时也成立，并认为粲企鹅可能在领头阶给出长程效应\cite{Bauer:2004tj}；BBNS 认为二者等价\cite{Beneke:2004bn}，BPRS 不同意\cite{Bauer:2005wb}，本书不作裁断。分歧集中在硬旁观者项能否按 $\as(\sqrt{m_b\Lqcd})$ 展开与粲企鹅的幂次计数（E-S 型陈述）；正因没有共识与公认证明，标签取 [P] 而非 [A]；但说“BBNS 只在 $\order{\as}$ 成立”则低估了现状。升级途径：包含这两点的公认全阶证明。

\begin{proposition}[幂次修正中的端点发散 \tagV]\label{prop:heavy-qcdf-endpoint}
在 $\order{\as}$ 按标准共线因子化形式计算 $\bar B\to M_1M_2$ 的扭度 3 手征增强硬旁观者项与弱湮灭项（$\Lqcd/m_b$ 幂次修正），即使取渐近分布振幅也得到发散的卷积，分别含
\begin{equation}\label{eq:heavy-XAXH}
X_H=\int_0^1\frac{\dd y}{1-y}\,\Phi_p(y),\qquad X_A=\int_0^1\frac{\dd y}{y},
\end{equation}
% VF-BQ-3
故这些修正的标准共线因子化不成立。参数化 $X_{A,H}=(1+\varrho_{A,H}\ee^{\ii\varphi_{A,H}})\ln\frac{m_B}{\Lambda_h}$（$\varrho_{A,H}\le1$，$\Lambda_h=0.5$\,GeV）是一个模型；例如对两个赝标介子、按 Beneke--Neubert 的规则 $\int_0^1\frac{\dd y}{y}\to X_A$，
\begin{equation}\label{eq:heavy-annih}
A_1^i\approx\pi\as\Big[18\Big(X_A-4+\frac{\pi^2}{3}\Big)+2r_\chi^2X_A^2\Big],\qquad A_3^f\approx12\pi\as\,r_\chi\big(2X_A^2-X_A\big).
\end{equation}
% VF-BQ-4, VF-BQ-5
（本命题是关于硬散射核与分布振幅之卷积的陈述。）
\uses{prop:heavy-qcdf-lp}
\end{proposition}

\begin{proof}[发散性的显式计算]
\cref{prop:heavy-qcdf-lp}中的 $T^{\rm II}$ 给出 $\int_0^1\frac{\dd y}{\bar y}[\Phi_{M_1}(y)+r_\chi^{M_1}\frac{\bar x}{x}\Phi_p(y)]$ 型积分\cite[式~(57)]{BenekeEtAl2001}：扭度 2 部分 $\int\frac{6y\bar y}{\bar y}=3$ 有限，$\Phi_p=1$ 时 $\int_0^{1-\delta}\frac{\dd y}{1-y}=\ln\frac1\delta$ 发散。发散来自旁观者作为软夸克进入 $M_1$ 的区域 $\bar y\to0$，由\cref{eq:heavy-hc-gluon}此时 $k^2\to\xi^2m_B^2\sim\Lqcd^2$；物理上它在 $\bar y\sim\Lqcd/m_b$ 处被截断，$X_H\sim\ln(m_b/\Lqcd)$，系数可能是复数\cite{BenekeEtAl2001}。弱湮灭构件 $A_1^i=\pi\as\int\!\!\int\dd x\,\dd y\{\Phi_{M_2}\Phi_{M_1}[\frac{1}{y(1-x\bar y)}+\frac{1}{\bar x^2y}]+r_\chi^{M_1}r_\chi^{M_2}\frac{2}{\bar xy}\}$\cite[式~(62)]{BenekeEtAl2001}：取渐近分布振幅，第一项有限，$36\int\!\!\int\frac{x\bar x\bar y}{1-x\bar y}=18(\frac{\pi^2}{3}-3)$；第二项 $18\int_0^1\frac{x\,\dd x}{1-x}$（$\frac{x}{1-x}=\frac1{1-x}-1$）与第三项都发散。把 $\int_0^1\frac{\dd y}{y}$、$\int_0^1\frac{\dd x}{1-x}$ 换成 $X_A$，三项为 $18(\frac{\pi^2}{3}-3)$、$18(X_A-1)$、$2r_\chi^2X_A^2$，相加即\cref{eq:heavy-annih}（$A_3^f$ 同理；脚本 \texttt{s2\_\allowbreak qcdf.py}）：“常数项”只在给发散积分规定数值之后才有意义。
\end{proof}

\paragraph{地位与缺口}
破坏（[V]），证据为显式计算。$X_{A,H}$ 只是对未知非微扰贡献的参数化：BBNS 称对所有湮灭项取同一个 $X_A$ 是“a crude model”\cite{BenekeEtAl2001}，Beneke 与 Neubert 称弱湮灭的处理“model-dependent”\cite{BenekeNeubert2003}；两文把默认值写作 $\ln(m_M/\Lambda_h)\approx2.4$，应为 $\ln(m_B/\Lambda_h)=2.36$（附录~Q 第~7~条）。把端点发散再因子化为有限的因子化定理是猜想性质的（[C]）：其他过程中已有逐例的成功，如经 $b$ 夸克圈的 $h\to\gamma\gamma$\cite{Liu:2019oav,Liu:2020wbn}、经轻夸克的 $gg\to h$\cite{Liu:2022ajh}、$\bar B\to X_s\gamma$ 的次领头 $O_8$--$O_8$ 贡献\cite{Hurth:2023paz}；对纯弱湮灭衰变，有文献补上一个此前遗漏的领头幂次硬共线贡献，消去了卷积中的快度发散\cite{Lu:2022kos}；色压低的 $\bar B^0\to D^0M^0$ 中弱湮灭与色压低振幅同阶，没有幂次压低\cite{Mantry:2003uz}。据我们截至 2026 年 10 月的文献检索（INSPIRE），尚未见对 $\bar B\to M_1M_2$ 的一般弱湮灭与扭度 3 旁观者项的完整再因子化（一般情形见第~21~章的\cref{prop:power-nlp-endpoint}，[P]，与\cref{sec:front-nlp}）。
